% Owner: Yogesh Modi. Literature Review (expanded from the Phase I text).
\section{Literature Review}
\label{sec:litreview}

\subsection{RISC and Pipelining}
Patterson and Ditzel argued that a reduced instruction set with simple,
uniform instructions allows a simpler and faster implementation than a complex
instruction set~\cite{yogeshPatterson1980}. The Stanford MIPS project showed
the same idea in silicon: a pipelined microprocessor in which the compiler and
the hardware cooperate to keep the pipeline busy~\cite{yogeshHennessy1982}.
Pipelining itself is older; Kogge describes it as overlapping the stages of
successive operations so that throughput is set by the slowest stage instead
of the whole operation~\cite{yogeshKogge1981}. The five-stage
IF--ID--EX--MEM--WB pipeline of Hennessy and Patterson became the standard
teaching model~\cite{yogeshHennessyQA,yogeshPattersonRV}, and Harris and Harris
give single-cycle, multicycle and pipelined RISC-V versions of it in
HDL~\cite{yogeshHarris2021}. Deeper pipelines raise the clock rate but pay more
cycles per hazard; Hartstein and Puzak show that this trade-off has an optimum
depth that depends on the workload~\cite{yogeshHartstein2002}. Our project
compares the two ends that matter for a teaching core: a single-cycle design
(CPI~$=1$, long clock period) and the classic five-stage pipeline (CPI~$>1$,
short clock period). It follows the textbook organisation without changes, so
its results can be checked against the textbook analysis.

\subsection{Hazards and Forwarding}
A pipeline must handle three kinds of hazard~\cite{yogeshHennessyQA}.
Structural hazards are avoided by giving instructions and data separate
memories, as in our Harvard design. For data hazards there are three main
approaches. The simplest is to stall the dependent instruction until the
producer has written the register file; this is correct but costs up to two
or three cycles for every dependent pair. \emph{Forwarding} (bypassing) instead
sends a result from the EX/MEM or MEM/WB register straight to the ALU input, so
most read-after-write dependencies cost no cycle; only a load followed by a
dependent instruction still needs one stall cycle, because the data arrives at
the end of MEM~\cite{yogeshPattersonRV}. Dynamic scheduling with Tomasulo's
algorithm~\cite{yogeshTomasulo1967} goes further and lets independent
instructions run ahead of a stalled one, but it needs reservation stations,
register renaming and a common data bus, which is far more hardware than an
in-order core. Our design uses forwarding plus a one-cycle load-use stall and a
write-first register file, which is the usual choice for a small in-order
pipeline: it removes most stall cycles at the cost of two 3-input multiplexers
and a comparator unit. The price is a longer EX path, since the forwarded value
passes through a multiplexer before the ALU.

\subsection{Branch Handling}
Control hazards arise because the next PC is known only after the branch is
resolved. Static schemes need no history: stall until the branch resolves,
always predict not-taken and flush on a taken branch, or predict by direction
(backward taken, forward not-taken)~\cite{yogeshHennessyQA}. Smith compared
these strategies with dynamic predictors that keep per-branch saturating
counters and found that a 2-bit counter table predicts much better than any
static rule~\cite{yogeshSmith1981}. Yeh and Patt's two-level adaptive predictor
adds branch history and reaches still higher accuracy~\cite{yogeshYeh1991}.
Dynamic predictors pay off in deep pipelines, where a misprediction costs many
cycles. In our five-stage design a taken branch or jump costs two cycles (it is
resolved in EX), and predict not-taken needs only the flush signals of the
hazard unit. A predictor is out of the scope of this project; its likely gain
can be estimated from the taken-branch counts measured in the performance
section.

\subsection{The RISC-V ISA and Open-Source Cores}
RISC-V defines a small base integer ISA (RV32I/RV64I) with optional standard
extensions~\cite{yogeshRVSpec}, and Asanovi\'c and Patterson argued for it as a
free and open ISA for academia and industry~\cite{yogeshAsanovic2014}. Its
fixed 32-bit encoding, with register fields always in the same place, makes it
easy to read registers in parallel with decoding, which suits a classic
pipeline. Open implementations cover a wide range of design points. The Rocket
Chip generator produces an in-order five-stage core with caches and virtual
memory~\cite{yogeshRocket2016}, and BOOM is a superscalar out-of-order core for
high performance~\cite{yogeshBoom2015}. At the small end, PicoRV32 is a
non-pipelined, multi-cycle core optimised for area and
Fmax~\cite{yogeshPicoRV32}, and the PULP cores trade pipeline depth against
energy for IoT devices: RI5CY has four stages with DSP
extensions~\cite{yogeshGautschi2017}, while Zero-riscy and Micro-riscy use two
stages to save area and power~\cite{yogeshSchiavone2017}. These cores show that
the number of stages is a design choice driven by the target (area, energy or
speed), not a fixed rule. Our core is far smaller than all of them (29
instructions, no caches, no CSRs). It is meant to be read and tested module by
module, and its single-cycle baseline makes the effect of pipelining
measurable on the same ISA, memories and test programs.

\subsection{Open-Source FPGA Synthesis Tools}
Synthesis results are usually obtained with vendor tools, which are closed and
tied to one device family. Yosys is a free Verilog synthesis suite that
supports most synthesizable Verilog and maps logic to FPGA primitives such as
LUTs, carry chains and block RAM~\cite{yogeshWolf2013}. Combined with the
placer and router nextpnr it forms a complete open flow from Verilog to
bitstream for commercial Lattice iCE40 and ECP5 devices, with timing-driven
placement and a static timing report~\cite{yogeshShah2019}. The academic VTR
flow also covers synthesis to routing and models many FPGA architectures, but
it targets research architectures rather than commercial
parts~\cite{yogeshMurray2020}. We use Yosys and nextpnr from the OSS CAD
Suite~\cite{yogeshOSSCAD} because they are free, scriptable and report
resource use and maximum frequency for real devices, which our course outline
asks for. One property of the target devices matters for our design: the block
RAM of iCE40 and ECP5 has a synchronous (clocked) read
port~\cite{yogeshICE40DS,yogeshECP5DS}, while our memories have a
combinational read. Whether a memory can use block RAM therefore depends on
whether the tool can absorb a register after (or before) the read, which is
discussed in the synthesis results.

\subsection{Reference Project and Gap Addressed}
We studied the open-source \emph{protocolx-riscv-pipelined-processor}
project~\cite{yogeshProtocolx}, which implements a pipelined RISC-V processor
in Verilog, while planning the module breakdown and the test approach. All
code in our repository was written independently by the team. Production cores
are too large to study line by line, and textbook datapaths are often given
only as diagrams with CPI estimates. This project fills the gap with a small,
fully tested five-stage design and a single-cycle baseline built from the same
modules. Both are measured with the same benchmarks and the same open-source
synthesis flow, so cycle counts and clock frequency are compared directly.
